%======================================================================
\section{Certification methods}\label{app:certification}
%======================================================================

This appendix describes the arithmetic, the validated special functions and the method for the exact members used in the certificates
of Section~\ref{sec:certified}; the conventions are stated in Section~\ref{ssec:cert-method}.

%----------------------------------------------------------------------
\subsection{Exact and ball arithmetic}
%----------------------------------------------------------------------

\begin{remark}[Ball arithmetic and rigorous solves]
\emph{Exact} computations use integers and rationals (FLINT \texttt{fmpz}, \texttt{fmpq} and their polynomial and matrix
types \cite{FLINT}). \emph{Interval} computations use Arb ball arithmetic \cite{Joh17}: a real number is represented by a
ball $[m\pm r]$ that is guaranteed to contain it, and every operation returns a ball containing all possible results.
Both are accessed through python-flint. Constants such as $\pi$ and Euler's constant $\gamma_E$ are balls. Linear systems are solved by Arb's
ball solver, which either returns balls containing the exact solution of every system in the input balls, and thereby certifies that
the exact matrix is invertible, or fails. The non-rigorous library mpmath \cite{mpmath} is used only in self-tests, never in a
certificate.
\end{remark}

Positivity of polynomials is certified by coefficient signs.

\begin{lemma}[Positive-coefficient certificates]\label{lem:posc}
\textup{(a)} \textup{(Descartes at $c$.)} If $p\in\RR[t]$ and every coefficient of $p(c+t)\in\RR[t]$ is non-negative with
positive constant term, then $p>0$ on $[c,\infty)$. If $p$ vanishes to exact order $m$ at $c$ and the coefficients of
$t^m,t^{m+1},\dots$ in $p(c+t)$ are positive, then $p>0$ on $(c,\infty)$.

\textup{(b)} A polynomial with non-negative coefficients in $v\in[0,\infty)^m$ is non-negative there, and positive at every
point at which some monomial with a positive coefficient is positive.

With ball coefficients, ``positive'' means that the ball lies in $(0,\infty)$.
\end{lemma}

\begin{proof}
Both statements follow from the fact that a polynomial with non-negative coefficients takes non-negative values
on the non-negative orthant, with the indicated strictness.
\end{proof}

In Proposition~\ref{prop:lowheight} the lemma is applied, after a M\"obius map of an interval onto $[0,\infty)$, to polynomials with
exact rational coefficients; Taylor shifts of ball polynomials, where needed, are computed by Horner's scheme at a ball containing $c$,
which encloses the exact shifted coefficients.

%----------------------------------------------------------------------
\subsection{Validated Bessel functions}
%----------------------------------------------------------------------

All certificates built on $\Xi^2$ (Propositions~\ref{thm:finiteJ} and~\ref{prop:exactcone}) evaluate $K_0$ and $K_1$ on positive real balls. We do not use the built-in Bessel
and confluent hypergeometric functions of Arb for this purpose: at moderate $z$ their enclosures can lose about $e^{2z}$
in relative accuracy. Instead the shipped module \nolinkurl{lib/besselk.py} uses two classical representations with
explicit error bounds.

\begin{lemma}[Validated evaluation of $K_0$ and $K_1$]\label{lem:K-validated}
Let $z>0$ and $\nu\in\{0,1\}$.
\textup{(a)} $K_0$ and $K_1$ are given by the convergent series \cite[10.31.2, 10.31.1]{DLMF}
\[
  \begin{aligned}
  K_0(z)&=-\bigl(\log\tfrac z2+\gamma_E\bigr)I_0(z)+\sum_{k\ge1}\frac{H_k\,q^k}{(k!)^2},\\
  K_1(z)&=\frac1z+\log\tfrac z2\,I_1(z)-\frac z4\sum_{k\ge0}\frac{(\psi(k+1)+\psi(k+2))\,q^k}{k!\,(k+1)!},
  \end{aligned}
\]
with $q=z^2/4$, $I_0(z)=\sum_kq^k/(k!)^2$, $I_1(z)=\frac z2\sum_kq^k/(k!(k+1)!)$ and $H_k$ the harmonic numbers.
\textup{(b)} For $\ell\ge\max(|\nu|-\frac12,1)$, $K_\nu(z)=\sqrt{\pi/(2z)}\,e^{-z}\bigl(\sum_{k<\ell}a_k(\nu)z^{-k}+R_\ell\bigr)$ with
$a_k(\nu)=\prod_{j=1}^k(4\nu^2-(2j-1)^2)/(k!\,8^k)$ and $|R_\ell|\le|a_\ell(\nu)|z^{-\ell}$ \cite[10.40.2 and
\S10.40(ii)]{DLMF}.
\textup{(c)} $K_0'=-K_1$ and $K_1'=-K_0-K_1/z$; hence $|K_0'|=K_1$ and $|K_1'|\le K_0+K_1/z$ on $(0,\infty)$.
\end{lemma}

\begin{proof}
(a) and (b) are the cited formulas, specialised to $\nu=0,1$; (c) follows from the recurrences for $K_\nu'$ and the
positivity of $K_0,K_1$ on $(0,\infty)$.
\end{proof}

\emph{Implementation.} For a ball $[m\pm r]\subset(0,\infty)$ the module evaluates $K_0,K_1$ at the exact midpoint $m$ and
adds $r$ times an upper bound for $|K_\nu'|$ on the ball, obtained from (c) and $(\log K_1)'\ge-1-1/w$. At the midpoint it
uses (b), summed until the remainder bound falls below $2^{-P-20}$ relative, when this is possible, and otherwise (a) at
working precision $P+2.885z+64$ bits, since the series for $K_\nu$ cancels the growing $I_\nu$-terms to a relative size
$e^{-2z}$. The series is truncated at an index $k\ge1$ with $4q\le(k+1)^2$. From there on every ratio of consecutive terms is at
most $\frac12$: the factor $q/(k+1)^2$ is at most $\frac14$ and decreases, and the coefficient ratios ($H_{k+1}/H_k$, and the ratios of the
values $\psi(k+1)+\psi(k+2)$, which are positive for $k\ge1$) are at most $2$. So each tail is at most twice its first neglected term,
hence at most the last computed term; twice the last computed term is added as a ball radius. Evaluating the series directly on a ball
would multiply the radius by about $e^{2z}$, which the midpoint evaluation avoids.

\emph{A trapezoid-rule evaluation.} The primary path of the certificates of Proposition~\ref{prop:exactcone} uses instead, for $z\ge1$,
the representations $K_0(z)=e^{-z}(2z)^{-1/2}\int_\RR f_0$ and $zK_1(z)=(2z)^{1/2}e^{-z}\int_\RR f_1$, where
$f_0(v)=e^{-v^2}(1+v^2/2z)^{-1/2}$ and $f_1(v)=v^2e^{-v^2}(1+v^2/2z)^{1/2}$ (\cite[10.32.8]{DLMF} with $t=1+v^2/z$), evaluated by the
trapezoid rule with step $h$ (\nolinkurl{lib/k01_trapezoid.py}). The $f_\nu$ are analytic in $|\Im v|<\sqrt{2z}$; for $0<a<\sqrt{2z}$,
\cite[Theorem~5.1]{TW14} bounds the discretisation error by $2M_\nu/(e^{2\pi a/h}-1)$, with $M_0=\sqrt\pi e^{a^2}(1-a^2/2z)^{-1/2}$ and
$M_1=\sqrt\pi e^{a^2}[(\frac12+a^2)+(\frac34+a^2+a^4)/4z]$ bounding $\int_\RR|f_\nu(x+iy)|\dd x$ for $|y|<a$ (both decrease in $z$ and are
taken at the lower end of the ball for $z$). The sum is truncated at $|k|\le n$, and the rest is bounded by geometric series, using
$f_0(v)\le e^{-v^2}$ and $f_1(v)\le3e^{-v^2/2}$. The second path of these certificates uses Lemma~\ref{lem:K-validated}.

%----------------------------------------------------------------------
\subsection{The exact members}\label{app:exact}
%----------------------------------------------------------------------

The certificates of Proposition~\ref{prop:exactcone} work with the exact polynomial $P_J$, without approximation or cushion. Their input
is the ball enclosure of $p^{(J)}$ of Proposition~\ref{thm:finiteJ}, and every bound below is computed from these balls, so it holds for
every coefficient vector in the box, in particular for the exact one. By Proposition~\ref{thm:bessel} and Lemma~\ref{lem:abel},
$\FT F_J^{(m)}(\xi)=(2\pi)^m2\pi\sqrt x\sum_Nd(N)(A_mK_0+B_mK_1)(2\pi Nx)$, where the ball polynomials $A_m,B_m$, of degree at most $D$, come
from exact recurrences for $(\theta+\frac12)^mP_J(-(\theta+\frac12)^2)\theta^2(\theta+1)^2$ acting on $\operatorname{span}\{z^kK_0,z^kK_1\}$.
The divisor sum is truncated with a rigorous tail bound from $d(N)\le2\sqrt N$, $0<K_0\le K_1\le K_{3/2}$ with
$K_{3/2}(w)=(\pi/2w)^{1/2}e^{-w}(1+1/w)$, and a geometric series in $N$. In the same way
$\lvert\FT F_J(\xi)\rvert\le\Upsilon(x):=4\pi(\lvert A_0\rvert+\lvert B_0\rvert)(z)\,e^{-z}(1+1/z)/(1-q)$, with $z=2\pi x$ and
$q=2^De^{-z}<\frac12$, where $\lvert A\rvert$ is the polynomial whose coefficients are the absolute values of those of $A$. Suprema of
$\lvert\FT F_J^{(m)}\rvert$ on an interval use that $\lvert A_m\rvert,\lvert B_m\rvert$ increase and $K_0,K_1$ decrease in $z$.

\emph{Nodes.} At a node $\xi_n$, $n\le n_J$, Proposition~\ref{prop:E} gives $\FT F_J(\xi_n)=\FT F_J'(\xi_n)=0$ exactly. So, with
$c_k=\FT F_J^{(k)}(\xi_n)/k!$ and $M\ge\sup\lvert\FT F_J^{(p+1)}\rvert$ on $[\xi_n-h,\xi_n+h]$, Taylor's formula gives
$\FT F_J(\xi_n+\delta)=\delta^2\bigl(\sum_{k=2}^pc_k\delta^{k-2}+R\delta^{p-1}\bigr)$ with $\lvert R\rvert\le M/(p+1)!$. If the least Bernstein
coefficient of the polynomial part on $[0,h]$ and on $[-h,0]$ exceeds $Mh^{p-1}/(p+1)!$, then $\FT F_J>0$ for $0<\lvert\delta\rvert\le h$,
and $\xi_n$ is a zero of order exactly two. The idea of subtracting the vanishing jets before dividing by $\delta^2$ comes from the
planar Cohn--Elkies certificate \cite{OAI26planar}; it makes a cushion unnecessary. The enclosures of $c_0$ and $c_1$ computed from the
coefficient balls are never used: the verifier only checks that they contain $0$, which detects gross input errors and nothing more. The
argument rests on Proposition~\ref{thm:finiteJ}, whose balls contain the exact coefficients, and on $c_0=c_1=0$, which holds only for the
exact $P_J$.

\emph{Gaps.} Between the node neighbourhoods, Taylor models of order $p\in[24,64]$ at exact dyadic centres give Bernstein lower bounds minus
Lagrange remainders on pieces with dyadic endpoints, which are bisected until the bound is positive. The neighbourhoods, rounded inward to
dyadic endpoints, and the pieces tile $[\xi_2,\xi(x_0)]$, or $[0,\xi(x_0)]$, where $\xi(x)=\log x/2\pi$; the tiling is checked in exact
arithmetic, and the item at $n=2$ starts at the exact $\xi_2$.

\emph{Far field.} With $Q_{P_J}$ and $g_{P_J}$ as in Lemma~\ref{lem:abel}, if every Taylor coefficient of $Q_{P_J}$ at $y_0=2\pi x_0$ is
positive, then $Q_{P_J}>0$ on $[y_0,\infty)$ (Lemma~\ref{lem:posc}(a)), so $g_{P_J}(z)>0$ for $z\ge y_0$, and every divisor term of
$\FT F_J(\xi)$ is positive when $x\ge x_0$. The verifier takes for $x_0$ the first of the points $n_J+k/20$, $k\ge1$, with this property; in
Table~\ref{tab:kappa-ladder}, $x_0-n_J\in[0.3,0.45]$.

\emph{The prime sum.} In $\Arch(F_J)=\frac1\pi\sum_{n>n_J}\Lambda(n)n^{-1/2}\FT F_J(\xi_n)$ the terms with $n\le n_{\max}$ ($400$ for
$J\le61$, $800$ for $J\ge110$) are enclosed one by one. Beyond, the terms are at most $\log n\cdot\Upsilon(n)/\pi$, and the ratio of
consecutive bounds is at most $((n+1)/n)^De^{-2\pi}\log(n+1)/\log n$, which decreases in $n$; so it is checked only at $n_1=n_{\max}+1$,
where it is some $q'<\frac12$, and the terms sum to at most $\log(n_1)\Upsilon(n_1)/((1-q')\pi)$. This gives the tail bounds of Proposition~\ref{prop:exactcone}, one for each $J$.

%----------------------------------------------------------------------
\subsection{The digamma weight}
%----------------------------------------------------------------------

The Gaussian--Laguerre certificates (Propositions~\ref{thm:kappaOPS}, \ref{prop:lowheight} and~\ref{thm:notcritical}) integrate against the
digamma weight $\Omega_{\mathfrak g}$ on a finite interval and bound the rest by an explicit tail, which uses the following estimate.

\begin{lemma}[The digamma function on vertical lines]\label{lem:digamma}
For $w>0$ and $\tau>0$, $\bigl|\Re\psi(w+i\tau)-\log\sqrt{w^2+\tau^2}\bigr|\le1/\tau$. In particular, if $0<w\le5\sqrt3$ and $t\ge10$, then
$0<\Re\psi(w+\frac{it}2)\le\log t+\frac2t$.
\end{lemma}

\begin{proof}
Let $f(x):=(x+w)/((x+w)^2+\tau^2)=\Re\,(x+w+i\tau)^{-1}$. The series for $\psi$ in Section~\ref{sec:setting} gives
$\Re\psi(w+i\tau)=\lim_{N\to\infty}\bigl(\log N-\sum_{n=0}^{N-1}f(n)\bigr)$, and $\int_0^Nf=\log\sqrt{(N+w)^2+\tau^2}-\log\sqrt{w^2+\tau^2}$. For each
$n$, $|f(n)-\int_n^{n+1}f|$ is at most the total variation of $f$ on $[n,n+1]$. On $[0,\infty)$, $f>0$ increases to its maximum $1/(2\tau)$,
attained where $x+w=\tau$ if that point lies in $[0,\infty)$, and then decreases to $0$; so its total variation there is at most $1/\tau$.
Hence $|\sum_{n<N}f(n)-\int_0^Nf|\le1/\tau$ for every $N$, and letting $N\to\infty$ gives the first claim. For the second take $\tau=t/2\ge5$:
then $\log\sqrt{w^2+\tau^2}\ge\log5>\frac15\ge\frac1\tau$, and $\log\sqrt{w^2+\tau^2}\le\log t$ because $w^2\le75\le\frac34t^2$.
\end{proof}

The tail estimates of \nolinkurl{general/verify_general.py} use the consequence $|\Re\psi(w+\frac{it}2)|\le\log t+2\le t$ for $t\ge10$.
The verifier enforces $T\ge10$ and, for every Gamma shift $k_j$, $0\le k_j$ and $(\frac12+k_j)^2\le300$, checked exactly; that is,
$w_j=\frac12(\frac12+k_j)\in[\frac14,5\sqrt3]$, where the lemma applies, and it rejects other inputs. The data used here have $k_j\in\{0,1\}$,
that is, $w=\frac14$ and $w=\frac34$. The restriction is needed: for fixed $t$, $\Re\psi(w+\frac{it}2)$ grows like $\log w$.

%----------------------------------------------------------------------
\subsection{Provenance}
%----------------------------------------------------------------------

Every input of a certificate is an exact object (integers, rationals, or decimal strings read as exact rationals or as balls) stored in a
file; for the cushioned certificates of \nolinkurl{kappa/} the file also stores the SHA-256 of a canonical string of the parameters, which the
verifier recomputes, and the certificates of the exact members check the SHA-256 of their input enclosure.
The scripts check the rounding direction of every printed decimal exactly. Each shipped log was produced by a fresh run of the shipped
script and records the command, the date, the software versions, the exit code and the wall time.
